\فصل{روش پیشنهادی و نتایج تجربی}

این فصل به تشریح روش پیشنهادی و ارائه نتایج تجربی اختصاص دارد.
در بخش نخست، طراحی الگوریتم ژنتیک بهبودیافته شامل
نمایش کروموزوم دوبخشی، عملگرهای ژنتیکی (تقاطع و جهش)،
سازوکار مقداردهی اولیه هوشمند، انتخاب تطبیقی عملگر
و جست‌وجوی محلی تشریح می‌شود.
در بخش دوم، تنظیمات آزمایشی شامل نمونه‌های معیار،
الگوریتم‌های مقایسه‌شده و پروتکل آزمایشی بیان می‌شود.
در بخش سوم، نتایج تجربی ارائه و تحلیل می‌شوند
و در نهایت تحلیل آماری و مطالعه حذفی انجام می‌شود.

\قسمت{روش پیشنهادی}

\زیرقسمت{نمای گرافیکی روش پیشنهادی}

برای ارائه تصویری کل‌نگر از اجزای روش، شکل~\رجوع{شکل:نمای-گرافیکی-روش}
گردش‌کار الگوریتم پیشنهادی را از مرحله تولید جمعیت اولیه تا استخراج جبهه نهایی
نمایش می‌دهد.
در این نما، ارتباط میان هسته \lr{NSGA-II}، سازوکار انتخاب تطبیقی عملگر،
ارزیابی مبتنی بر شبیه‌ساز و جست‌وجوی محلی به‌صورت یکپارچه مشخص شده است.

\شروع{شکل}[H]
\centerimg{proposed_method_pipeline.pdf}{14cm}
\شرح{نمای گرافیکی روش پیشنهادی شامل مقداردهی اولیه هوشمند، \lr{NSGA-II}، انتخاب تطبیقی عملگر، ارزیابی شبیه‌سازی‌محور و جست‌وجوی محلی}
\برچسب{شکل:نمای-گرافیکی-روش}
\پایان{شکل}

مطابق شکل~\رجوع{شکل:نمای-گرافیکی-روش}، جریان حل به‌صورت یک حلقه بازخوردی طراحی شده است:
نتیجه ارزیابی هر نسل مستقیماً به سازوکار انتخاب عملگر بازمی‌گردد و در نسل بعد اثر می‌گذارد.
این پیوستگی بین «ارزیابی»، «به‌روزرسانی پاداش» و «تولید نسل بعد» دلیل اصلی بهبود تدریجی کیفیت جبهه پارتو در طول اجراست.


\زیرقسمت{نمایش کروموزوم}

در این پایان‌نامه از نمایش دوبخشی برای کروموزوم استفاده می‌شود:

\شروع{شمارش}
\فقره \مهم{بخش ترتیب عملیات}\پاورقی{Operation Sequence}:
یک جایگشت از شماره کارها به طول $n \times m$.
هر بار ظاهر شدن شماره کار $j$، به عملیات بعدی آن کار اشاره دارد.
این نمایش تضمین می‌کند که محدودیت تقدم عملیات هر کار همواره رعایت شود~\cite{bierwirth1995generalized}.
\فقره \مهم{بخش نگاشت ماشین}\پاورقی{Machine Map}:
یک بردار به طول $m$ که هر ماشین منطقی را به یکی از گره‌های رایانشی
(لبه، مه یا ابر) نگاشت می‌کند.
\پایان{شمارش}


\زیرقسمت{عملگرهای ژنتیکی}

چهار عملگر تقاطع\پاورقی{Crossover} و چهار عملگر جهش\پاورقی{Mutation}
پیاده‌سازی شده‌اند:

\مهم{عملگرهای تقاطع}:
\شروع{فقرات}
\فقره \lr{JOX}\پاورقی{Job Order Crossover}: تقاطع مبتنی بر ترتیب کارها
\فقره \lr{PPX}\پاورقی{Precedence Preserving Crossover}: تقاطع حفظ‌کننده تقدم
\فقره تقاطع حفظ‌کننده کار\پاورقی{Job-Preserving Crossover}
\فقره تقاطع یکنواخت\پاورقی{Uniform Crossover}
\پایان{فقرات}

\مهم{عملگرهای جهش}:
\شروع{فقرات}
\فقره جهش جابه‌جایی\پاورقی{Swap Mutation}: جابه‌جایی دو ژن
\فقره جهش درج\پاورقی{Insert Mutation}: حذف و درج ژن در موقعیت جدید
\فقره جهش درهم‌ریزی\پاورقی{Scramble Mutation}: درهم‌ریزی یک زیربازه
\فقره جهش تخصیص مجدد\پاورقی{Reassign Map}: تغییر نگاشت ماشین به گره
\پایان{فقرات}


\زیرقسمت{مقداردهی اولیه هوشمند}

در الگوریتم ژنتیک بهبودیافته، جمعیت اولیه با ترکیب دو منبع ساخته می‌شود:

\شروع{شمارش}
\فقره \مهم{بذرهای ابتکاری}: راه‌حل‌هایی که توسط قواعد \lr{SPT} و \lr{MWR} تولید شده‌اند
به‌عنوان نقاط شروع با کیفیت بالا به جمعیت اولیه اضافه می‌شوند.
\فقره \مهم{راه‌حل‌های تصادفی}: بخش باقی‌مانده جمعیت با جایگشت‌های تصادفی و نگاشت‌های تصادفی ماشین تکمیل می‌شود تا تنوع حفظ شود.
\پایان{شمارش}


\زیرقسمت{انتخاب تطبیقی عملگر}

به‌جای استفاده از یک عملگر تقاطع و جهش ثابت،
سازوکار \lr{UCB1}~\cite{auer2002finite} برای انتخاب پویای عملگرها استفاده شده است.
در هر نسل، عملگری انتخاب می‌شود که بالاترین امتیاز \lr{UCB1} را دارد.
پاداش هر عملگر بر اساس بهبود ابرحجم فرزندان نسبت به والدین محاسبه می‌شود.
این سازوکار تعادل بین بهره‌برداری\پاورقی{Exploitation}
و اکتشاف\پاورقی{Exploration} را برقرار می‌کند.


\زیرقسمت{جست‌وجوی محلی}

در الگوریتم بهبودیافته، با احتمال ۳۰ درصد،
جست‌وجوی محلی روی فرزندان تولیدشده اعمال می‌شود.
جست‌وجوی محلی شامل بررسی همسایه‌های مجاور
(از طریق جابه‌جایی و درج عملیات)
و پذیرش بهبود در صورت غلبه پارتو است.


\زیرقسمت{شبه‌کد الگوریتم پیشنهادی}

برای یکپارچه‌سازی مراحل اجرایی، روند کلی الگوریتم پیشنهادی
در الگوریتم~\رجوع{الگوریتم:بهبودیافته} خلاصه شده است.

\شروع{الگوریتم}{الگوریتم ژنتیک چندهدفه بهبودیافته}
\برچسب{الگوریتم:بهبودیافته}
\ورودی نمونه \lr{JSSP}، پیکربندی زیرساخت، اندازه جمعیت $N$، تعداد نسل $G$
\خروجی جبهه راه‌حل‌های غیرمغلوب

\دستور جمعیت اولیه را با ترکیب بذرهای \lr{SPT}/\lr{MWR} و افراد تصادفی تولید کن
\دستور جمعیت اولیه را ارزیابی کن و رتبه غلبه و فاصله ازدحام را محاسبه کن
\تاوقتی{نسل فعلی از $G$ کمتر است}
\دستور والدین را با تورنمنت دودویی انتخاب کن
\دستور عملگرهای تقاطع و جهش را با \lr{UCB1} انتخاب کن
\دستور با اعمال عملگرهای انتخاب‌شده، فرزندان را تولید کن
\دستور با احتمال از پیش‌تعیین‌شده، جست‌وجوی محلی را روی فرزندان اعمال کن
\دستور فرزندان را با شبیه‌ساز ارزیابی کن و پاداش عملگرها را به‌روزرسانی کن
\دستور جمعیت والدین و فرزندان را ادغام کرده و بهترین $N$ فرد را نگه‌دار
\پایان‌تاوقتی
\دستور جبهه غیرمغلوب نهایی را برگردان
\پایان{الگوریتم}


\قسمت{تنظیمات آزمایشی}

\زیرقسمت{پیاده‌سازی و جریان ابزارها}

برای شفاف‌سازی فرآیند پیاده‌سازی، جریان ابزارهای مورد استفاده در این پژوهش
در شکل~\رجوع{شکل:جریان-ابزارها} نمایش داده شده است.
این جریان شامل بارگذاری نمونه‌های معیار، کدگذاری کروموزوم، اجرای هسته تکاملی،
ارزیابی سناریوها در شبیه‌ساز \lr{YAFS} و در نهایت محاسبه شاخص‌ها و آزمون‌های آماری است.
این سازمان‌دهی، ردیابی فرایند تولید نتایج و بازتولید آزمایش‌ها را تسهیل می‌کند.

\شروع{شکل}[H]
\centerimg{implementation_toolflow.pdf}{14cm}
\شرح{جریان ابزارها در پیاده‌سازی: از ورودی نمونه‌ها تا تولید شاخص‌ها، تحلیل آماری و مصنوعات خروجی}
\برچسب{شکل:جریان-ابزارها}
\پایان{شکل}

همان‌گونه که در شکل~\رجوع{شکل:جریان-ابزارها} مشاهده می‌شود، خروجی هر مرحله ورودی مرحله بعدی است
و هیچ گسست دستی در خط لوله وجود ندارد؛ ازاین‌رو ردیابی خطا و بازتولید نتایج ساده‌تر می‌شود.
این سازمان‌دهی همچنین تضمین می‌کند که شاخص‌های گزارش‌شده (مانند \lr{HV} و \lr{IGD})
مستقیماً از همان اجراهایی محاسبه شوند که نمودارها و تحلیل‌های آماری بر مبنای آن‌ها تولید شده‌اند.


\زیرقسمت{نمونه‌های معیار}

آزمایش‌ها روی ۱۵ نمونه مسئله از مجموعه داده‌های استاندارد~\cite{weise2019jsspinstances}
انجام شده‌اند.
در جدول~\رجوع{جدول:نمونه‌ها} مشخصات نمونه‌های مورد استفاده فهرست شده است.

\شروع{لوح}[ht]
\تنظیم‌ازوسط
\شرح{نمونه‌های معیار مورد استفاده در آزمایش‌ها}

\شروع{جدول}{|l|c|c|c|}
\خط‌پر
\سیاه نمونه & \سیاه ابعاد ($n \times m$) & \سیاه خانواده & \سیاه بهترین مقدار شناخته‌شدهٔ زمان تکمیل کل \\
\خط‌پر \خط‌پر
\lr{ft06} & $6 \times 6$ & \lr{FT} & 55 \\
\lr{la01} & $10 \times 5$ & \lr{LA} & 666 \\
\lr{la06} & $15 \times 5$ & \lr{LA} & 926 \\
\lr{la11} & $20 \times 5$ & \lr{LA} & 1222 \\
\lr{la16} & $10 \times 10$ & \lr{LA} & 945 \\
\lr{la21} & $15 \times 10$ & \lr{LA} & 1046 \\
\lr{la26} & $20 \times 10$ & \lr{LA} & 1218 \\
\lr{orb05} & $10 \times 10$ & \lr{ORB} & 887 \\
\lr{abz5} & $10 \times 10$ & \lr{ABZ} & 1234 \\
\lr{abz7} & $20 \times 15$ & \lr{ABZ} & 656 \\
\lr{ta01} & $15 \times 15$ & \lr{TA} & 1231 \\
\lr{ta11} & $20 \times 15$ & \lr{TA} & 1357 \\
\lr{ta21} & $20 \times 20$ & \lr{TA} & 1642 \\
\lr{swv01} & $20 \times 10$ & \lr{SWV} & 1407 \\
\lr{yn1} & $20 \times 20$ & \lr{YN} & 884 \\
\خط‌پر
\پایان{جدول}

\برچسب{جدول:نمونه‌ها}
\پایان{لوح}


\زیرقسمت{الگوریتم‌های مقایسه‌شده}

شش الگوریتم مقایسه شده‌اند:

\شروع{شمارش}
\فقره \lr{heuristic\_spt}: قاعده کوتاه‌ترین زمان پردازش (خط پایه)
\فقره \lr{heuristic\_mwr}: قاعده بیشترین کار باقی‌مانده (خط پایه)
\فقره \lr{plain\_ga}: الگوریتم \lr{NSGA-II} ساده (جمعیت ۱۰۰، ۳۰۰ نسل)
\فقره \lr{enhanced\_ga}: الگوریتم \lr{NSGA-II} بهبودیافته (جمعیت ۱۰۰، ۴۰۰ نسل)
\فقره \lr{ablation\_no\_aos}: بدون انتخاب تطبیقی عملگر
\فقره \lr{ablation\_no\_local\_search}: بدون جست‌وجوی محلی
\پایان{شمارش}

هر الگوریتم در قالب ۱۰ اجرای مستقل (با بذرهای تصادفی متفاوت) اجرا شده است.


\زیرقسمت{پارامترهای اجرایی}

پارامترهای اصلی اجرای الگوریتم ژنتیک بهبودیافته در جدول~\رجوع{جدول:پارامترها} (پیوست) ارائه شده است.
در نسخهٔ پایه \lr{plain\_ga}، اندازه جمعیت ۱۰۰ و تعداد نسل‌ها ۳۰۰ در نظر گرفته شد؛
در حالی‌که برای الگوریتم بهبودیافته، اندازه جمعیت ۱۰۰ و تعداد نسل‌ها ۴۰۰ انتخاب شد.
دو آزمایش حذفی نیز با همین تنظیمات اجرا شدند، با این تفاوت که مؤلفهٔ متناظر در هر حالت غیرفعال شد.

\زیرقسمت{آزمون‌های آماری}

برای اعتبارسنجی آماری نتایج از آزمون‌های زیر استفاده شده است:

\شروع{فقرات}
\فقره \مهم{آزمون فریدمن}~\cite{friedman1937use}:
آزمون ناپارامتری برای مقایسه چند الگوریتم روی چند نمونه مسئله.
\فقره \مهم{آزمون ویلکاکسون}~\cite{wilcoxon1945individual}:
آزمون ناپارامتری دوتایی برای مقایسه هر جفت الگوریتم.
مقادیر $p$ با تصحیح هولم~\cite{holm1979simple} تعدیل شده‌اند.
\فقره \مهم{اندازه اثر}:
از شاخص‌های \lr{A12} وارگا-دلانی~\cite{vargha2000critique}
و دلتای کلیف\پاورقی{Cliff's delta} استفاده شده است.
\پایان{فقرات}


\قسمت{نتایج تجربی}

\زیرقسمت{جبهه‌های پارتو}

شکل~\رجوع{شکل:پارتو-شبکه} نمای کلی جبهه‌های پارتو را برای تمام نمونه‌ها نشان می‌دهد.
الگوریتم ژنتیک بهبودیافته در اکثر نمونه‌ها جبهه پارتو بهتری تولید کرده است.

\شروع{شکل}[H]
\centerimg{pareto_grid.pdf}{14cm}
\شرح{نمای کلی جبهه‌های پارتو (زمان تکمیل کل در برابر مصرف انرژی) برای تمام نمونه‌های معیار}
\برچسب{شکل:پارتو-شبکه}
\پایان{شکل}

همان‌گونه که در شکل~\رجوع{شکل:پارتو-شبکه} مشاهده می‌شود، جبهه‌های حاصل از
\lr{enhanced\_ga} در بیشتر نمونه‌ها نسبت به \lr{plain\_ga} و روش‌های ابتکاری
به ناحیه غالب نزدیک‌ترند. تحلیل کمی نتایج نیز نشان می‌دهد در مقایسه چهار روش اصلی
(\lr{enhanced\_ga}، \lr{plain\_ga}، \lr{heuristic\_mwr} و \lr{heuristic\_spt})،
\lr{enhanced\_ga} در هر ۱۵ نمونه هم بیشترین ابرحجم و هم کمترین \lr{IGD} را کسب کرده است.

\زیرقسمت{شاخص ابرحجم}

شکل~\رجوع{شکل:ابرحجم-جعبه‌ای} توزیع مقادیر ابرحجم را برای هر نمونه نشان می‌دهد.
الگوریتم‌های ژنتیکی به‌طور معنادار از قواعد ابتکاری بهتر عمل کرده‌اند.
الگوریتم بهبودیافته در نمونه‌های بزرگ‌تر برتری بیشتری نشان می‌دهد.

\شروع{شکل}[H]
\centerimg{hv_boxplot_per_instance.pdf}{14cm}
\شرح{توزیع ابرحجم برای هر نمونه مسئله (بیشتر بهتر است)}
\برچسب{شکل:ابرحجم-جعبه‌ای}
\پایان{شکل}

بر اساس شکل~\رجوع{شکل:ابرحجم-جعبه‌ای} و میانگین نتایج روی ۱۵ نمونه و ۱۰ اجرای مستقل،
\lr{enhanced\_ga} نسبت به \lr{plain\_ga} به‌طور متوسط حدود $0.65\%$ ابرحجم بالاتری دارد.
در مقایسه با \lr{heuristic\_mwr}، این بهبود به بیش از ۱۶ برابر می‌رسد و نسبت به
\lr{heuristic\_spt} چند مرتبه بزرگی بیشتر است. همچنین نسبت برد/باخت در ابرحجم
برای \lr{enhanced\_ga} در برابر \lr{plain\_ga} برابر $15/0$ بوده است.

\زیرقسمت{شاخص فاصله نسلی معکوس}

شکل~\رجوع{شکل:ای‌جی‌دی-جعبه‌ای} توزیع مقادیر \lr{IGD} را نشان می‌دهد.
مقادیر کمتر نشان‌دهنده نزدیکی بیشتر به جبهه مرجع تجربی است.

\شروع{شکل}[H]
\centerimg{igd_boxplot_per_instance.pdf}{14cm}
\شرح{توزیع فاصله نسلی معکوس برای هر نمونه مسئله (کمتر بهتر است)}
\برچسب{شکل:ای‌جی‌دی-جعبه‌ای}
\پایان{شکل}

تحلیل شکل~\رجوع{شکل:ای‌جی‌دی-جعبه‌ای} نشان می‌دهد \lr{enhanced\_ga}
به‌صورت پایدار مقدار \lr{IGD} را کاهش داده است. به‌طور متوسط،
\lr{IGD} این روش نسبت به \lr{plain\_ga} حدود $45.8\%$ کمتر است؛
این کاهش در مقایسه با \lr{heuristic\_mwr} و \lr{heuristic\_spt}
به‌ترتیب حدود $99.5\%$ و $99.8\%$ است.

\زیرقسمت{همگرایی}

شکل~\رجوع{شکل:همگرایی} منحنی‌های همگرایی ابرحجم را در طول نسل‌ها نشان می‌دهد.
نوار سایه‌دار نشان‌دهنده یک انحراف معیار در ۱۰ بذر تصادفی است.
الگوریتم بهبودیافته با سرعت بیشتری همگرا می‌شود.

\شروع{شکل}[H]
\centerimg{convergence_global.pdf}{12cm}
\شرح{منحنی همگرایی ابرحجم میانگین روی تمام نمونه‌ها}
\برچسب{شکل:همگرایی}
\پایان{شکل}

همان‌طور که در شکل~\رجوع{شکل:همگرایی} دیده می‌شود، سرعت همگرایی نسخه بهبودیافته بالاتر است.
بر اساس داده‌های همگرایی، \lr{enhanced\_ga} حدود نسل ۱۵ به ۹۵٪ کیفیت نهایی خود می‌رسد،
در حالی که این مقدار برای \lr{plain\_ga} حدود نسل ۳۰ است. همچنین در نسل‌های ۵۰، ۱۰۰ و ۲۰۰،
ابرحجم میانگین \lr{enhanced\_ga} تقریباً $2.7$ برابر نسخه پایه گزارش شده است.


\زیرقسمت{نتایج بهترین زمان تکمیل کل}

شکل~\رجوع{شکل:زمان‌تکمیل-جعبه‌ای} بهترین زمان تکمیل کل یافته‌شده توسط هر الگوریتم را نشان می‌دهد.

\شروع{شکل}[H]
\centerimg{boxplot_best_makespan.pdf}{14cm}
\شرح{توزیع بهترین زمان تکمیل کل برای هر نمونه}
\برچسب{شکل:زمان‌تکمیل-جعبه‌ای}
\پایان{شکل}

مطابق شکل~\رجوع{شکل:زمان‌تکمیل-جعبه‌ای}، \lr{enhanced\_ga}
در ۱۵ نمونه، کمترین زمان تکمیل کل را در میان چهار روش اصلی به‌دست آورده است.
به‌صورت کمی، میانگین زمان تکمیل کل این روش نسبت به \lr{plain\_ga} حدود $5.3\%$
و نسبت به \lr{heuristic\_mwr} و \lr{heuristic\_spt} به‌ترتیب حدود $82.5\%$ و $93.2\%$
کمتر است. این نتیجه نشان می‌دهد بهبود کیفی جبهه پارتو فقط در شاخص‌های چندهدفه نیست،
بلکه در معیار کلاسیک زمان تکمیل نیز حفظ شده است.


\زیرقسمت{نمودار گانت}

شکل~\رجوع{شکل:گانت} نمودار گانت بهترین برنامه زمان‌بندی یافته‌شده
توسط الگوریتم بهبودیافته را نشان می‌دهد.

\شروع{شکل}[H]
\centerimg{gantt_representative.pdf}{14cm}
\شرح{نمودار گانت بهترین برنامه زمان‌بندی از الگوریتم ژنتیک بهبودیافته}
\برچسب{شکل:گانت}
\پایان{شکل}

در تفسیر شکل~\رجوع{شکل:گانت} مشاهده می‌شود که ترتیب عملیات به‌گونه‌ای تنظیم شده
که هم‌پوشانی مؤثر بین ماشین‌ها افزایش یابد و زمان‌های بی‌کاری طولانی کاهش پیدا کند.
چگالی مناسب اجرای عملیات در ماشین‌های مختلف، یکی از عوامل اصلی کاهش
$C_{\max}$ در نسخه بهبودیافته است.


\قسمت{تحلیل آماری}

\زیرقسمت{آزمون فریدمن}

برای هر شاخص، آزمون فریدمن بر روی نتایج الگوریتم‌ها در مجموعهٔ نمونه‌های معیار اعمال شد.
در تمامی موارد، مقدار $p$ به‌طور معناداری کمتر از ۰.۰۵ به دست آمد
که وجود تفاوت آماری معنادار میان الگوریتم‌ها را تأیید می‌کند.

\زیرقسمت{نقشه حرارتی اندازه اثر}

شکل~\رجوع{شکل:اهمیت-ابرحجم} نقشه حرارتی دلتای کلیف را
برای مقایسه الگوریتم بهبودیافته با سایر الگوریتم‌ها نشان می‌دهد.
حاشیه ضخیم نشان‌دهنده تفاوت آماری معنادار ($p_{\text{adj}} < 0.05$) است.

\شروع{شکل}[H]
\centerimg{significance_heatmap_hypervolume.pdf}{14cm}
\شرح{نقشه حرارتی دلتای کلیف برای شاخص ابرحجم. حاشیه ضخیم: $p_{\text{adj}} < 0.05$}
\برچسب{شکل:اهمیت-ابرحجم}
\پایان{شکل}

بررسی دقیق شکل~\رجوع{شکل:اهمیت-ابرحجم} نشان می‌دهد در مقایسه با
\lr{heuristic\_mwr} و \lr{heuristic\_spt}، برتری \lr{enhanced\_ga}
در هر ۱۵ نمونه هم مثبت و هم از نظر آماری معنادار است
($p_{\text{adj}}<0.05$ و $\delta>0$). در مقایسه با \lr{plain\_ga} نیز
اندازه اثر در تمام نمونه‌ها مثبت بوده و در ۶ نمونه معنادار گزارش شده است؛
بنابراین برتری نسخه بهبودیافته صرفاً بصری نبوده و پشتوانه آماری دارد.


\زیرقسمت{رتبه‌بندی الگوریتم‌ها}

شکل~\رجوع{شکل:رتبه-حرارتی} نقشه حرارتی رتبه میانگین ابرحجم را نشان می‌دهد.
رتبه ۱ بهترین عملکرد را نشان می‌دهد.

\شروع{شکل}[H]
\centerimg{rank_heatmap.pdf}{12cm}
\شرح{نقشه حرارتی رتبه میانگین ابرحجم برای هر نمونه و هر الگوریتم (۱ = بهترین)}
\برچسب{شکل:رتبه-حرارتی}
\پایان{شکل}

مطابق شکل~\رجوع{شکل:رتبه-حرارتی}، \lr{enhanced\_ga} در هر ۱۵ نمونه رتبه ۱ را
برای شاخص ابرحجم کسب کرده است و \lr{plain\_ga} عمدتاً رتبه ۲ را دارد.
این الگو با نتایج آزمون‌های آماری هم‌راستا است و پایداری عملکرد روش پیشنهادی
را در مجموعه متنوعی از ابعاد مسئله نشان می‌دهد.


\قسمت{مطالعه حذفی}

برای سنجش سهم مؤلفه‌های بهبودی بررسی‌شده، دو حالت حذفی اجرا شده‌اند.
شکل~\رجوع{شکل:حذفی} اندازه اثر حذف هر مؤلفه را نشان می‌دهد.

\شروع{شکل}[H]
\centerimg{ablation_effect_sizes.pdf}{12cm}
\شرح{اندازه اثر حذف هر مؤلفه بهبودی (مقادیر منفی = کاهش عملکرد پس از حذف)}
\برچسب{شکل:حذفی}
\پایان{شکل}

نتایج مطالعه حذفی نشان می‌دهد:
\شروع{فقرات}
\فقره حذف جست‌وجوی محلی بیشترین افت را در شاخص ابرحجم ایجاد می‌کند.
\فقره حذف انتخاب تطبیقی عملگر نیز اثر منفی متوسطی دارد.
\پایان{فقرات}

به‌صورت کمی و بر اساس شکل~\رجوع{شکل:حذفی}، حذف جست‌وجوی محلی نسبت به
\lr{enhanced\_ga} باعث افت متوسط حدود $0.43\%$ در ابرحجم و بدترشدن
\lr{IGD} به میزان حدود $31.8\%$ می‌شود. حذف انتخاب تطبیقی عملگر نیز
به افت متوسط حدود $0.20\%$ در ابرحجم و افزایش حدود $23.0\%$ در \lr{IGD}
منجر شده است. این نتایج نشان می‌دهد هر دو مؤلفه در کیفیت نهایی نقش دارند
و سهم جست‌وجوی محلی برجسته‌تر است.


\قسمت{مقایسه بصری کلی}

شکل~\رجوع{شکل:رادار} مقایسه راداری الگوریتم‌ها را روی پنج شاخص نشان می‌دهد.

\شروع{شکل}[H]
\centerimg{radar_comparison.pdf}{10cm}
\شرح{مقایسه راداری الگوریتم‌ها روی شاخص‌های مختلف}
\برچسب{شکل:رادار}
\پایان{شکل}

برداشت کمی از شکل~\رجوع{شکل:رادار} آن است که روش پیشنهادی در اکثر محورها
به ناحیه بهینه نزدیک‌تر است. میانگین رتبه این روش در پنج شاخص اصلی
حدود $1.13$ گزارش شده، در حالی‌که این مقدار برای \lr{plain\_ga}
حدود $1.87$ است. این اختلاف نشان می‌دهد بهبودها به یک شاخص خاص محدود نیستند
و اثر سیستماتیک بر کیفیت کلی دارند.

شکل~\رجوع{شکل:برد-باخت} نسبت برد/تساوی/باخت الگوریتم بهبودیافته در برابر سایر الگوریتم‌ها را نشان می‌دهد.

\شروع{شکل}[H]
\centerimg{win_tie_loss.pdf}{12cm}
\شرح{نسبت برد/تساوی/باخت الگوریتم بهبودیافته بر اساس شاخص ابرحجم}
\برچسب{شکل:برد-باخت}
\پایان{شکل}

بر اساس شکل~\رجوع{شکل:برد-باخت}، در شاخص ابرحجم نسبت
برد/تساوی/باخت \lr{enhanced\_ga} در برابر \lr{plain\_ga}،
\lr{heuristic\_mwr} و \lr{heuristic\_spt} به‌ترتیب
$15/0/0$، $15/0/0$ و $15/0/0$ است.
در مقایسه با \lr{ablation\_no\_aos} این نسبت $9/3/3$ و در برابر
\lr{ablation\_no\_local\_search} برابر $15/0/0$ به‌دست آمده است
که بار دیگر نقش مثبت مؤلفه‌های بهبودی را تأیید می‌کند.


\قسمت{جمع‌بندی فصل}

در این فصل، ابتدا اجزای الگوریتم پیشنهادی شامل نمایش کروموزوم، عملگرهای ژنتیکی،
مقداردهی اولیه هوشمند، انتخاب تطبیقی عملگر و جست‌وجوی محلی به‌صورت منسجم تشریح شد.
سپس تنظیمات کامل آزمایش، شاخص‌های ارزیابی و پروتکل مقایسه ارائه شد.
در نهایت، نتایج تجربی و تحلیل‌های آماری نشان دادند که الگوریتم بهبودیافته
در اغلب نمونه‌ها عملکرد بهتری نسبت به روش‌های پایه ارائه می‌دهد.
